% !TeX root = booklet-01.tex
\bibitem{intro}
\lr{Introduction}، از مجموعه اسلایدهای \lr{CMOS VLSI Design, Fourth Edition}، ۲۷ اسلاید.
طبق اسلاید ۲، یادداشت‌های اصلی اثر \lr{David Money Harris} (۲۰۱۱) است و
اقتباس‌کننده با حروف \lr{MS} معرفی شده است. شماره اسلاید و صفحه فایل یکسان‌اند.
داده‌های تاریخی در متن با زمینه زمانی مرجع تفسیر شده‌اند.
\bibitem{mos}
\lr{MOSFET Transistors}، از مجموعه اسلایدهای \lr{CMOS VLSI Design, Fourth Edition}،
برگرفته از یادداشت‌های \lr{David Money Harris} (۲۰۱۱)، ۲۵ اسلاید.
از ساختار و ناحیه کاری تا مثال جریان، مدل ظرفیت و ولتاژ تغذیه استفاده شده است.
در اسلاید ۱۵، برخی روابط در لایه متن وجود دارند، اما در تصویر صفحه جای خالی دارند.
روابط مربوط با کتاب تطبیق داده شده‌اند.
\bibitem{current}
\begin{latin}
Samuel J. Ling, William Moebs, and Jeff Sanny, \textit{University Physics Volume 2},
OpenStax, 2016, Section 9.1,
\href{https://openstax.org/books/university-physics-volume-2/pages/9-1-electrical-current}{Electrical Current}.
\end{latin}
\nopagebreak[4]
مرجع تکمیلی تعریف جریان و جهت قراردادی در فصل پیش‌نیاز.
\bibitem{resistors}
\begin{latin}
Samuel J. Ling, William Moebs, and Jeff Sanny, \textit{University Physics Volume 2},
OpenStax, 2016, Section 10.2,
\href{https://openstax.org/books/university-physics-volume-2/pages/10-2-resistors-in-series-and-parallel}{Resistors in Series and Parallel}.
\end{latin}
\nopagebreak[4]
مرجع تکمیلی مقاومت و اتصال‌های سری و موازی.
\bibitem{kirchhoff}
\begin{latin}
Samuel J. Ling, William Moebs, and Jeff Sanny, \textit{University Physics Volume 2},
OpenStax, 2016, Section 10.3,
\href{https://openstax.org/books/university-physics-volume-2/pages/10-3-kirchhoffs-rules}{Kirchhoff's Rules}.
\end{latin}
\nopagebreak[4]
مرجع تکمیلی قوانین گره و حلقه.
\bibitem{capacitors}
\begin{latin}
Samuel J. Ling, William Moebs, and Jeff Sanny, \textit{University Physics Volume 2},
OpenStax, 2016, Section 8.1,
\href{https://openstax.org/books/university-physics-volume-2/pages/8-1-capacitors-and-capacitance}{Capacitors and Capacitance}.
\end{latin}
\nopagebreak[4]
مرجع تکمیلی خازن و تعریف ظرفیت.
